% Generated draft: requires mathematical and source review.
\documentclass[11pt]{article}
\usepackage[margin=25mm]{geometry}
\usepackage{fontspec}
\setmainfont{Latin Modern Roman}
\usepackage{amsmath,amssymb,mathtools,unicode-math}
\setmathfont{Latin Modern Math}
\usepackage{enumitem,needspace,xurl,hyperref}
\setlength{\parindent}{0pt}
\setlength{\parskip}{5pt}
\setlength{\emergencystretch}{3em}
\newcommand{\N}{\mathbb{N}}
\newcommand{\Z}{\mathbb{Z}}
\newcommand{\Q}{\mathbb{Q}}
\newcommand{\R}{\mathbb{R}}
\newcommand{\C}{\mathbb{C}}
\newcommand{\F}{\mathbb{F}}
\newcommand{\A}{\mathbb{A}}
\newcommand{\PP}{\mathbb{P}}
\newcommand{\E}{\mathbb{E}}
\newcommand{\eps}{\varepsilon}
\newcommand{\dd}{\,\mathrm{d}}
\begin{document}
\section*{UnsolvedMath ID: 30001907}
\textbf{Source problem code:} catalogue entry
\subsection{Definitions and mathematical statement}
Let \(E\) be Morava \(E\)-theory of height \(n\), and let \(q\) be odd, as stipulated in the supplied record. Write \(S^q\) for the sphere, \(\Sigma^k\) for suspension by \(k\), and \(\widetilde E_q(S^q)\) for reduced \(E\)-homology in degree \(q\). Retain the source notation \(D_{2^k}(\operatorname{Id})\), \(E_q^{\wedge}\), \(C[k]_{-q}^{\vee}\), \(\Phi_n\), and \(\beta_k\). The supplied identification is
\[
E_q^{\wedge}\!\left(\Sigma^kD_{2^k}(\operatorname{Id})(S^q)\right)
\cong C[k]_{-q}^{\vee}.
\]
The question is whether, under this identification, the induced map \((\Phi_n\beta_k)_*\) on \(E\)-homology is the corresponding differential in the Koszul resolution of \(\widetilde E_q(S^q)\).

This is a conditional comparison question requiring the source conventions. The record does not specify the prime, the admissible ranges of \(n,q,k\) beyond oddness of \(q\), the completion and duality conventions, the definition of \(C[k]\), the source and target of \(\beta_k\), or the algebra and grading conventions for the Koszul resolution. In particular, it does not supply the target identification or the sign and indexing conventions needed to state an explicit equality between homomorphisms. These qualifications remain unresolved; the displayed identification alone does not determine that equality.
\subsection{Short English statement}
For Morava \(E\)-theory of height \(n\) and odd \(q\), the problem asks whether the map induced by \(\Phi_n\beta_k\), under the supplied identification, agrees with a Koszul differential for \(\widetilde E_q(S^q)\). Several definitions and map conventions needed for a fully specified equality are missing. The supplied partially solved label is a dated curation assessment, not an independently verified conclusion.
\subsection{Sources}
\begin{itemize}
\item Supplied local UnsolvedMath registry record, permanent ID \(30001907\), title \emph{Koszul Differentials in Bousfield--Kuhn Homology}, exact source problem code \texttt{catalogue entry}; collection: Oberwolfach Reports Open Problems. No individual proposer is named; the supplied proposed year is 2011.
\item Supplied report citation: \emph{Homotopy Theory (2012)}, Oberwolfach Reports, pp.\ 2623--2667, DOI \texttt{10.4171/owr/2011/46}. Original OWR report: \url{https://doi.org/10.4171/owr/2011/46}. The record describes these reports as documenting mathematical workshops and their open-problem discussions.
\item Supplied literature reference: Behrens--Rezk, \emph{Bousfield--Kuhn and topological Andre--Quillen cohomology}, \url{https://arxiv.org/abs/1712.03045}.
\item Literature assessment date: \texttt{2026-08-21}. This is supplied dated curation, not independently verified. It labels the problem partially solved and reports substantial Bousfield--Kuhn/Koszul comparison results, including an odd-degree Koszul description and sphere cases, while stating that the extracted conjecture is not established in its stated generality. The supplied background calls this dated literature triage and requests independent verification of time-sensitive status and cited claims.
\item The registry labels the statement \texttt{corrected\_verified} and reports a check against the official Oberwolfach report preserving the hypotheses, notation, and open target. This verification claim is supplied metadata; no independent source inspection is asserted here.
\end{itemize}
\end{document}
